\documentclass{memoireuqamV2}

% --- PRÉAMBULE ---
\usepackage{times}
\usepackage[utf8]{inputenc}
\usepackage{lmodern}
\usepackage[T1]{fontenc}
\usepackage[french]{babel}
\renewcommand{\tablename}{Tableau}
\usepackage{url}
\usepackage{xcolor}

\makeatletter
\g@addto@macro{\UrlBreaks}{\do\/\do\-\do\.}
\makeatother

\usepackage{graphicx}
\usepackage{amsmath}
\usepackage{amsfonts}
\usepackage{amssymb}
\usepackage{booktabs}
\usepackage{array} 
\usepackage{ragged2e} 
\usepackage{multirow}
\usepackage{natbib}
\usepackage{multibib}
\newcites{refII}{Références II}
\usepackage{microtype}
\usepackage{caption}
\usepackage{subcaption}
\usepackage{enumitem}
\usepackage{setspace}
\usepackage{float}  
\usepackage{tocloft}

\usepackage[breaklinks=true]{hyperref} 

% --- CONFIGURATIONS ---
\hypersetup{
    colorlinks=true,
    linkcolor=black,
    filecolor=black,
    urlcolor=black,
    citecolor=black,
}

% --- MÉTA-DONNÉES DU DOCUMENT ---
\title{Examen général de doctorat en informatique cognitive (DIC-9401) \\ 
 \vspace{6em}
Le Léviathan algorithmique : pouvoir, opacité et responsabilité à l'ère de l'intelligence artificielle \\ \vspace{6em}}
\author{Dominique Loyer}
\uqamthese
\matiere{informatique cognitive}


% --- DÉBUT DU DOCUMENT ---
\begin{document}

\onehalfspacing
\maketitle
\clearpage

% --- RÉSUMÉ ---
\input{resume.tex}

% --- TABLE DES MATIÈRES ---
\onecolumn
\small
\tableofcontents
\newpage

% --- CORPS DU RAPPORT ---
\setlength{\parskip}{1em}

\input{introduction.tex}
\input{chapitre1.tex}
\input{chapitre2.tex}
\input{chapitre3.tex}
\input{chapitre4.tex}
\input{chapitre5.tex}
\input{chapitre6.tex}
\input{chapitre7.tex}
\input{conclusion.tex}

% --- BIBLIOGRAPHIE ---
\bibliographystyle{apalike-uqam}
\renewcommand{\refname}{Références citées}
\bibliography{refI}

\subsection*{Références non-citées mais utilisées}
\renewcommand{\refname}{Références non-citées}
\bibliographystyle{apalike-uqam}
\bibliography{refII}

% --- ANNEXES ---
\clearpage
\appendix
\input{annexes.tex}


\end{document}
